\section{Newton 法用曲率校正方向}
\label{sec:09-Convex-Optimization/06-Optimization-Algorithms/02}

梯度下降知道往哪边走是下坡，但不知道每个方向的下坡有多陡。它把函数在当前点附近近似成一个倾斜的平面——只有一阶信息。平面无法告诉你曲率：这一脚踩下去，地面是平缓延伸还是突然陡降？

你拿着梯度走一步，步长按最陡方向统一设定，结果在平缓方向浪费步数，在陡峭方向越界震荡。上节那个 \(f(x,y)=\frac12x^2+50y^2\) 的狭长碗就是明证——条件数 100 就够让朴素梯度下降步履蹒跚，之字形锯齿般折返几百步才能摸到谷底。你也许会用预条件或更好的步长策略缓解，但这些办法都绕着一个根本问题走：它们没有直接把曲率信息编码进方向本身。就像你闭着一只眼看路——知道大致走向，但不知道每步踩下去地有多软多硬。

Newton 法的思路直接：既然一阶信息不够，那就保留二阶。在当前点 \(x\) 附近做二阶 Taylor 展开：

\[
f(x+p)
\approx
f(x)+\nabla f(x)^{\trans}p
+\frac12p^{\trans}\nabla^2f(x)p.
\]

这是一个二次模型。一阶项告诉你哪边下坡；二阶项告诉你下坡的过程中坡面如何弯曲——曲率在各方向的分布被 Hessian 完整编码。如果 Hessian \(\nabla^2f(x)\) 在当前位置正定，这个二次函数有一个唯一的最小值点，满足

\[
\nabla^2f(x)p=-\nabla f(x).
\]

解出方向 \(p_{\mathrm N}\) 后，更新

\[
x_+=x+\alpha p_{\mathrm N}.
\]

Newton 步做的事情不只是``梯度除以二阶导数''的高维搬运——那个口诀丢了全部几何。它做的是：先用 Hessian 把各方向的曲率差异校正掉——相当于伸缩并旋转坐标系，让局部地形在变换后的坐标中接近一个标准圆——再在校正后的坐标系中找到二次模型的极小点。方向和步长在一次求解中同时确定：方向被 Hessian 扭转（不再垂直等高线，而是直接瞄向二次模型的底部），步长被 Hessian 的逆自动缩放（曲率大的方向走短些，曲率小的方向走长些）。

这就是 Newton 法比梯度下降贵得多的原因：每一步不仅要求梯度，还要求 Hessian 并解一个线性系统。付出的代价换回了什么？换回了方向的质量——一个有曲率感知的方向，能让你在一步中跨过梯度下降需要几十上百步才能走完的距离。

\subsection{二次函数为什么一步到达}

如果你面对的本就是一个二次函数：

\[
f(x)=\frac12x^{\trans}Qx+b^{\trans}x+c,
\qquad Q\succ0.
\]

梯度为 \(Qx+b\)，Hessian 恒为 \(Q\)。Newton 方程给出

\[
Qp=-(Qx+b),
\]

所以

\[
x+p=-Q^{-1}b=x^\star.
\]

二次模型就是原函数本身，精确算术中一步到达最优点。不管 \(Q\) 的条件数有多大——\(10^2\)、\(10^6\)、\(10^{10}\)——Newton 法都在一步内完成梯度下降可能需要数千步才能完成的事。不是快了一点，是快了一个量级跃迁。

这解释了两件事。第一，上节让梯度下降反复折返的狭长二次碗——对 Newton 法完全不是问题：Hessian 把两个方向的曲率差异一并吸收进线性系统，一步就把各向异性地形变成了各向同性。条件数在 Newton 法面前不是问题，因为 Hessian 逆矩阵已经精确地对每个方向做了反比例缩放。第二，更一般的光滑函数只在靠近最优点时才近似二次——所以 Newton 法最强的速度兑现于局部。远离最优时，高阶效应主导地形，二次模型不再可靠，Newton 步可能会指向完全错误的方向，甚至方向所指向的预测下降和实际函数值变化完全无关。

\subsection{局部二次收敛需要条件}

靠近最优点时，若 Hessian 正定且变化足够平滑（精确条件通常表述为 Hessian 在最优点的某个邻域内 Lipschitz 连续），Newton 步（取 \(\alpha=1\)）的误差满足

\[
\norm{ x_{k+1}-x^\star}
\le
C\norm{ x_k-x^\star}^2.
\]

误差每次``平方自己''。梯度下降的收敛因子是常数（几何收敛），Newton 法的速率却是把一个越来越小的数字不断平方：当前误差是 \(10^{-2}\)，下一步约 \(C\cdot10^{-4}\)，再下一步 \(C^3\cdot10^{-8}\)。正确数字的位数翻倍增长——从 \(10^{-2}\) 到 \(10^{-16}\) 只需三四步，而梯度下降靠几何收缩要达到同样精度可能需要数百步。

来看一个具体对比。从 \(x_0\) 出发，\(f(x)=x^4\)（光滑但 Hessian 在零点退化，这只是一个近似的量级说明），几步后的误差量级如下。梯度下降（固定步长）：\(10^{-1}, 9\times10^{-2}, 8.1\times10^{-2},\ldots\)，每一步误差只缩减约一成。Newton 法：\(10^{-1}, 10^{-2}, 10^{-4}, 10^{-8}, 10^{-16}\)，四步就从一位精度跳到机器精度。这就是几何收敛和二次收敛之间的实际差距——不是一个温和的常数倍数，而是一种收敛方式的质变。

但必须重申：这一结果是带余项的 Taylor 展开加上 Hessian 连续性导出的——它是一个局部定理，不是全局保证。离最优点太远时，二次模型可能完全失准：函数的实际下降可能远小于模型预测，甚至方向可能指向上升（函数值增大）。即使原问题凸，如果 Hessian 接近奇异或条件极差，Newton 方程本身也接近病态——Hessian 的微小扰动被求逆操作放大到主导方向，求解出的方向可能更多反映数值噪声而非真实地形。非凸问题中 Hessian 甚至不定——具有负特征值——这时二次模型在负曲率方向上无下界，求解出的 \(p\) 指向一个鞍点而非极小点，沿着 \(p\) 走函数值可能反而上升。

因此实际使用的算法几乎都会``阻尼''：解出 Newton 方向后，不直接走 \(\alpha=1\)，而是用回溯线搜索挑一个 \(\alpha\le 1\)。远离最优时，线搜索可能把 \(\alpha\) 压到 \(0.1\)、\(0.01\) 甚至更小——这时 Newton 方向的价值并不比梯度下降大太多。但一旦进入二次模型可靠的局部区域，线搜索连续若干步都接受 \(\alpha=1\)，算法自动从``保命模式''切换到二次收敛的``冲刺模式''。

阻尼 Newton 法是一种``低速巡航后高速冲刺''的策略。前期的低速不是可选的——在二次模型不可靠的区域硬走 \(\alpha=1\) 会导致发散。线搜索在前期扮演的不是性能优化器，而是安全阀。

\subsection{实际计算通过线性系统得到 Newton 步}

Newton 方向应通过求解线性系统

\[
H p=-g
\]

得到。你不需要——也绝不应该——显式算出 \(H^{-1}\)。矩阵求逆在数值上较不稳定，其浮点运算次数也大约是求解线性系统的三倍。永远解 \(Hp=-g\)，永远不算 \(H^{-1}\)。对于一个 \(n=100\) 的问题，\(H\) 是 \(100\times100\) 矩阵——求逆和求解的效率差距不大。但对于 \(n=10000\) 的问题（机器学习中常见），\(H\) 是 \(10000\times10000\)，求逆需要 \(O(n^3)\approx 10^{12}\) 次浮点运算，求解（通过 Cholesky 或共轭梯度）可以降到 \(O(n^2)\) 到 \(O(n^3/3)\) 的范围。对于 \(n=10^6\)：Hessian 稠密时存储就需要约 8TB——物理上不可行。

在大规模机器学习问题中，一个 Transformer 的参数有数亿维——存储 Hessian 需要 \(10^{16}\) 个浮点数。

这时有两类折中。截断 Newton 和 Hessian-free 方法只近似求解 Newton 方程——通过共轭梯度（CG）迭代到较低精度。关键技巧：CG 每次迭代只需计算 Hessian-vector product \(Hv\)，而 \(Hv\) 可以通过自动微分在一次额外的反向传播中算出，无需形成 \(H\) 本身。通常在共轭梯度迭代 5-50 轮后就截断，每轮成本约等于一次梯度求值。总成本是梯度下降每步的若干倍，但方向质量远优于朴素梯度下降。

拟 Newton 方法（BFGS、L-BFGS）走另一条路：从梯度变化序列构造 Hessian 逆的低秩近似。核心公式是 BFGS 更新：
\[
H_{k+1}^{-1} \approx \left(I - \frac{s_k y_k^{\trans}}{y_k^{\trans}s_k}\right) H_k^{-1} \left(I - \frac{y_k s_k^{\trans}}{y_k^{\trans}s_k}\right) + \frac{s_k s_k^{\trans}}{y_k^{\trans}s_k},
\]
其中 \(s_k=x_{k+1}-x_k\)，\(y_k=\nabla f(x_{k+1})-\nabla f(x_k)\)。L-BFGS 不存储整个矩阵，只保留最近 \(m\) 对 \((s_k,y_k)\)（典型 \(m=10\) 到 \(50\)），每次迭代用两层循环执行拟 Newton 方向的矩阵-向量乘法，成本 \(O(mn)\)。存储为 \(O(mn)\) 而非 \(O(n^2)\)。

比较算法时，不要问``二阶方法每步是不是更快''——它每一步更贵，且贵得多。应该问：达到所需精度的总时间是多少？对大规模稀疏问题，L-BFGS 每步约是梯度下降的 2-5 倍成本，但步数可能是梯度下降的 1/10 到 1/100——总时间显著更少。对稠密小规模问题，完整 Newton 每步成本是梯度下降的 \(n\) 倍量级，但步数极少——需要具体问题具体权衡。

\subsection{Newton decrement 衡量局部模型还剩多少}

Hessian 正定时，可以定义一个量来衡量在当前二次模型中还能降多少：

\[
\lambda(x)^2 = \nabla f(x)^{\trans} \nabla^2f(x)^{-1} \nabla f(x)
= -\nabla f(x)^{\trans}p_{\mathrm N}.
\]

把它除以 2，\(\lambda(x)^2/2\) 恰好是当前二次模型预测的最优下降量：把 Newton 方向代入二次模型的函数值预测，得到的减少量就是它。如果二次模型完美描述了原函数，这一步之后函数值就下降 \(\lambda(x)^2/2\)。

Newton decrement 有两重用途。第一，作为停止判据：如果 \(\lambda(x)^2/2\) 已经小于预设的容差（如 \(10^{-8}\)），说明在当前二次模型的精度范围内，已经探测不到任何显著的下降空间——再降需要借助三阶及更高阶效应，那些早已超出 Newton 法的感知范畴，继续迭代无意义。第二，对一类特殊的凸函数——自协调函数（self-concordant functions），其三阶导数被 Hessian 自身按一种不依赖于坐标系的方式控制——Newton decrement 能给函数值次优度一个精确的可计算上界。在后续内点法章节中，它将被用作核心的停止准则和步长决策依据。

再次提醒：计算时不要求逆。解出 \(p_{\mathrm N}\) 后，直接计算 \(-\nabla f(x)^{\trans}p_{\mathrm N}\) 即得 decrement。如果已在求解 Newton 方程，这就是一个极便宜的额外向量乘积。

\subsection{正则化与信赖域}

Hessian 不定或接近奇异时，Newton 方程的解毫无意义：不定矩阵的``逆''会把负曲率方向的能量放大到任意的程度，方向指向的可能是一个上升方向。补救办法是在 Hessian 上加一个正则项：

\[
(H+\mu I)p=-g.
\]

\(\mu>0\) 把全体特征值向上平移，保证正定性。调控 \(\mu\) 的效果：\(\mu\to 0\) 时恢复原始 Newton 步；\(\mu\) 很大时，\(H+\mu I\approx\mu I\)，方向趋近于负梯度方向，步长自动缩小为原来的 \(1/\mu\)。这个连续插值连接了 Newton 方向和梯度方向——\(\mu\) 充当了``曲率信任度''的旋钮：曲率信息可靠时用小的 \(\mu\)（近乎 Newton），不可靠时用大的 \(\mu\)（近乎梯度下降）。

另一种思路是信赖域方法：不修改 Hessian，而是限制步长在局部模型可信的半径内：

\[
\min_{\norm{ p}\le\Delta}
f(x)+\nabla f(x)^{\trans}p+\frac12p^{\trans}\nabla^2f(x)p.
\]

\(\Delta\) 的大小依据模型预测和实际函数值下降的匹配程度动态调整。如果实际下降接近甚至超过模型预测——模型是好的——放大 \(\Delta\)。如果实际下降远小于模型预测——模型已失准——缩小 \(\Delta\)，让下一步更保守。Levenberg--Marquardt 方法可以看作这两种思路的等价体：在某个意义下，增大 \(\mu\) 和缩小 \(\Delta\) 产生相同的限制效果。

在严格的凸优化主线中，Hessian 半正定已经比非凸的 Hessian 不定情况温和得多；但数值误差、边界附近的弱曲率、和有限精度导致的近似秩亏损，仍然要求正则化。理论上的曲率性质和计算出的矩阵条件数是两回事——前者可能温和，后者可能因特征尺度和单位选择而极其恶劣。两者需要同时关注，尤其在用迭代法（如共轭梯度）近似求解 Newton 方程时：CG 的收敛速度取决于计算矩阵的条件数，而不是理论 Hessian 的条件数。
